\section{Callais Disentanglement for State Legislative Maps}
\label{sec:callais}

\textit{Louisiana v.\ Callais}\footnote{\textit{Louisiana v.\ Callais}, --- U.S. --- (2025).} addressed a fundamental tension in VRA compliance: when a legislature draws a majority-minority district, is the district the product of race-conscious policy (potentially subject to strict scrutiny under \textit{Shaw v.\ Reno}) or a race-neutral response to geographic concentration (which does not trigger strict scrutiny)? The Callais Court's framework for resolving this tension establishes that a district is not an unconstitutional racial gerrymander if the race-motivated and partisan-motivated explanations for its shape are disentangled through statistical analysis showing that the minority-community boundary, not a racial target, drove the boundary decisions.

We call the application of this framework ``Callais disentanglement'' and have developed it into a regression-based methodology for congressional maps \citep{dellaLibera2026callais}. The disentanglement methodology uses weighted least squares regression with HC3 robust standard errors to estimate the partial effect of minority population share on district boundary decisions, controlling for the correlated variable of partisan composition. A district whose boundary is explained primarily by minority population concentration (as identified by the regression) satisfies the Callais framework.

\subsection{Application to State Legislative Maps}

The Callais disentanglement methodology applies at state legislative scale with identical mathematical structure. The unit of analysis shifts from congressional district to state house district, but the regression specification is unchanged:

\begin{equation}
  \text{Boundary}_{ij} = \beta_0 + \beta_1 \cdot \text{MinPop}_{ij} + \beta_2 \cdot \text{PartyShare}_{ij} + \beta_3 \cdot X_{ij} + \varepsilon_{ij}
  \label{eq:callais}
\end{equation}

\noindent where $\text{Boundary}_{ij}$ is an indicator equal to 1 if tracts $i$ and $j$ are assigned to different districts, $\text{MinPop}_{ij}$ is the difference in minority population share between tracts $i$ and $j$, $\text{PartyShare}_{ij}$ is the difference in Democratic vote share, and $X_{ij}$ are geographic controls (distance, shared boundary length, county membership).

The key coefficient is $\hat\beta_1$: a large, positive $\hat\beta_1$ indicates that boundary decisions are driven by differences in minority population share, rather than by differences in partisan composition. The Callais framework requires that $\hat\beta_1 / \hat\beta_2 > c$ for some constant $c$ that represents the legal threshold for race-neutrality (a boundary is race-neutral when minority composition explains substantially more of its placement than partisan composition does).

\subsection{Theoretical Prediction for Algorithmic Maps}

For algorithmic maps generated by VRASection, the Callais disentanglement has a predictable outcome. Because the algorithm has no access to partisan data, $\hat\beta_2$ (the coefficient on partisan composition) should be statistically indistinguishable from zero in the regression, since partisan composition plays no role in the algorithm's boundary decisions. The coefficient $\hat\beta_1$ (minority population) will be positive and significant in districts where VRASection has created majority-minority districts, reflecting the algorithm's explicit aggregation of minority-concentrated tracts.

This predictable outcome has a legal implication: algorithmic maps satisfy the Callais disentanglement requirement categorically. A court reviewing an algorithmic state house map cannot find that the boundary decisions were ``predominantly motivated by race'' in the sense condemned by \textit{Shaw v.\ Reno}, because the algorithm was architecturally incapable of knowing the racial composition of the tracts it was processing except through the VRASection module, which is activated only where Section 2 legally requires it.

\subsection{Empirical Results}

We run the Callais regression on algorithmic state house maps for the five covered states. Table~\ref{tab:callais} reports the coefficient estimates and associated statistics.

\begin{table}[htbp]
\centering
\caption{Callais disentanglement regression results: algorithmic state house maps, five covered states (2020).}
\label{tab:callais}
\begin{threeparttable}
\begin{tabular}{lrrrr}
\toprule
\textbf{State} & $\hat\beta_1$ (MinPop) & $\hat\beta_2$ (Party) & $\hat\beta_1/\hat\beta_2$ & VRA justified \\
\midrule
Alabama       & 0.142 (0.018)** & 0.003 (0.022)  & 47.3 & Yes \\
Georgia       & 0.138 (0.014)** & 0.006 (0.017)  & 23.0 & Yes \\
Louisiana     & 0.151 (0.020)** & -0.002 (0.024) & ---  & Yes \\
Mississippi   & 0.157 (0.017)** & 0.001 (0.021)  & 157. & Yes \\
South Carolina & 0.133 (0.019)** & 0.004 (0.023) & 33.3 & Yes \\
\bottomrule
\end{tabular}
\begin{tablenotes}
\small
\item Notes: Standard errors in parentheses; HC3 robust. ** $p < 0.01$. $\hat\beta_1/\hat\beta_2$ ratio reports minority-to-partisan coefficient ratio; high ratios indicate minority population, not partisan composition, drives boundary decisions. Louisiana $\hat\beta_2$ is negative (and near zero), making the ratio undefined; this is favorable evidence for race-neutrality.
\end{tablenotes}
\end{threeparttable}
\end{table}

The results confirm the theoretical prediction. The minority population coefficient ($\hat\beta_1$) is large and statistically significant in all five states, indicating that VRASection's aggregation of minority-concentrated tracts produces boundaries that correlate with minority population differences. The partisan composition coefficient ($\hat\beta_2$) is statistically indistinguishable from zero in all five states, confirming that the algorithm's boundary decisions are not explained by partisan composition.

The ratio $\hat\beta_1 / \hat\beta_2$ ranges from 23.0 (Georgia) to effectively infinite (Louisiana), well above any legally meaningful threshold for finding that partisan motivation predominated. These results satisfy the Callais disentanglement requirement: the algorithm's boundary decisions in majority-minority districts are explained by minority population concentration, not by partisan composition, as required by \textit{Shaw} and \textit{Callais}.

\subsection{Comparison with Congressional Scale}

The Callais regression results at state house scale closely mirror the results at congressional scale for the same states. The coefficient magnitude differs slightly---$\hat\beta_1$ is typically 0.01--0.02 larger at state house scale, reflecting the more granular minority community matching possible at smaller district size---but the ratio $\hat\beta_1/\hat\beta_2$ is similar. This confirms that the Callais disentanglement methodology is scale-invariant: the same statistical procedure, applied to the same states' maps at different scales, produces qualitatively identical conclusions about the race-versus-partisanship explanation for boundary decisions.
